\section{Binary Search: the Lower Bound (open)}\label{sec:lb}

\subsection{The program}
\begin{lstlisting}
method main(a: array<int, int>, n: int, x: int) returns (lo: int)
  requires 0 <= n
  requires sorted(a, 0, n)
  ensures 0 <= lo <= n
  ensures forall k: int :: 0 <= k < lo ==> a[k] < x
  ensures forall k: int :: lo <= k < n ==> x <= a[k]
{
  var hi: int; var mid: int;
  lo := 0; hi := n;
  while lo < hi {
    mid := (lo + hi) / 2;
    if a[mid] < x { lo := mid + 1; } else { hi := mid; }
  }
}
\end{lstlisting}
The first position of $a[0..n)$ whose element is at least $x$, for a sorted array. The example
\texttt{bsearch.dfy} of \texttt{examples/dafny/teq} verifies a binary search with the stronger
precondition of partition points for all indices; here the precondition is $\sorted(a, 0, n)$, so the
comparisons $a[m] < x$ are meaningful for $0 \le m < n$ only.

\subsection{The trace equation system}
\begin{align}
  K &= \cnt(0, 1, n') \qquad lo = \cur(L) \qquad hi = \cur(H) \qquad hi \le lo \label{eq:lb0}\\
  L &= \tr{0}{\bl{K}, \bl{H}}{lo, k, hi \mapsto \upd{G}{(hi + lo)/2 + 1}\ \ \upd{\neg G}{lo}} \label{eq:lb1}\\
  H &= \tr{n}{\bl{K}, \bl{L}}{hi, k, lo \mapsto \upd{G}{hi}\ \ \upd{\neg G}{(hi + lo)/2}} \label{eq:lb2}\\
  &\qquad\text{where } G \text{ abbreviates } a[(hi + lo)/2] < x \nonumber\\
  &\al{\bl{K}, \bl{L}, \bl{H}}{k, lo, hi \mapsto lo < hi} \label{eq:lb3}
\end{align}
As in the merge, the two variables are traces over each other, with a data-dependent guard that reads the
array at a position computed from both.

\subsection{The attempt}
The intended lemmas name the result by an abstract function $\mathsf{lb}(a, n, x)$ with the facts
$0 \le \mathsf{lb}(a, n, x) \le n$, $(a[k] < x) \Leftrightarrow \mathsf{true}$ for $0 \le k < \mathsf{lb}(a, n, x)$,
and $(x \le a[k]) \Leftrightarrow \mathsf{true}$ for $\mathsf{lb}(a, n, x) \le k < n$ (all under $\sorted(a, 0, n)$),
which close the postcondition once $lo = \mathsf{lb}(a, n, x)$ is known. The summary of the loop in the style of
the merge (Section~\ref{sec:merge}): the guard is the higher-order $G(lo, hi)$ pinned by
$\forall lo, hi.\ G(lo, hi) \Leftrightarrow a[(hi + lo)/2] < x$, the trace of $hi$ is tied by an open trace variable
$Hv = \tr{n}{\bl{\cnt(0, 1, nn)}, \bl{Lv}}{\ldots}$ with $hs = \bl{Hv}$, and the left-hand side is
$\cur(\tr{0}{\bl{\cnt(0, 1, nn)}, hs}{lo, k, hi \mapsto \upd{G(lo, hi)}{(hi + lo)/2 + 1}\ \ \upd{\neg G(lo, hi)}{lo}})$,
rewritten to $\mathsf{lb}(a, n, x)$. What the lemma must also know is that the loop has terminated, the exit
$hi \le lo$ of \eqref{eq:lb0}: without it the current value of $L$ is not the lower bound.

Stating the exit as a side condition $\cur(Hv) \le \cur(Lv)$ makes the run exceed the time limit: to decide
the condition, the simplifier rewrites $\cur(Lv)$, which is the left-hand side of the very lemma, whose
conditions decide it again (finding 20 of \texttt{NOTES.md}). Stating it as a comparison $h \le l$ of two
open variables, matched against the hypothesis $hi \le lo$, makes the lemma fire, and the four facts about
$\mathsf{lb}$ then close the three postconditions: the goal is closed, and a wrong postcondition leaves a
goal. But $h$ and $l$ are not tied to the traces, so the rule is not a sound lemma (in this goal its
conclusion holds whichever comparison the condition matched, since the exit does hold in the final state).
Every sound tie fails for a reason of the engine, located with probes: the instance of a side condition is
simplified before it is decided or matched, so $l = \cur(Lv)$ with $Lv$ bound re-enters the lemma like the
first form, whether closed or open. With $Lv$ open, $l = \cur(Lv)$ binds both from the hypothesis $lo = \cur(L)$
without re-entrance; but then the left-hand side, matched against $\cur(L)$ through the definition of the
trace variable $L$, turns out to have left its dependency variables $nn$ and $hs$ unbound (a probe fired with
"variables of a rule not instantiated"), and the trace condition that would bind them,
$Lv = \tr{0}{\bl{\cnt(0, 1, nn)}, hs}{\ldots}$, no longer matches once $Lv$ is bound, since its instance is
simplified with $L$ replaced by its definition.

\subsection{The idea, and what is missing}
Binary search is the paper's acknowledged weakness: the loop computes a bound by narrowing, and only a
partial description of the traces (monotone $lo$, monotone $hi$, $lo \le \mathsf{lb} \le hi$ throughout) is
needed, not their values. \texttt{bsearch.dfy} handles it by abstracting the comparisons with the array
into comparisons of indices with partition points (lemmas $(a[m] < x) \Leftrightarrow m < p$ without side
conditions), by the named monotone traces $\mathsf{incs}$ and $\mathsf{decs}$ whose bounds are their
facts, and by a cut after the loop. With $\sorted$ instead of partition points, the comparison lemma needs
$0 \le m < n$, which in the context of the updates follows from the invariant $0 \le lo \le hi \le n$,
the thing to be established; the summary lemma of this section bypasses the invariant by stating the
result directly, and founders on tying the exit to the traces soundly. What is missing is in the engine: a
rule should not be applied while the instances of its own side conditions are simplified, and a left-hand
side matched through a trace variable should bind all its variables; with either, the sound form of the
lemma closes the goal.
